\scanpage{157}
\mspage{V.29}
% archive sheet number: 156

6. \uline{Cas d'un point de non lissité \inserted{dans $X_0$} maximal}. \deleted{\unsure{de $X_0$}}

(6.1) Soit $x\in X_0$, et \deleted{suppo} reprenons les notations du § 4. Nous supposons que l'on a pour tt $i$
\[
\tag*{(6.1.1)}
{\unsure{\underline{\Phi}|\underline{\varphi}}}^i_{\unsure{\nu|v}}|V=0,\quad\text{i.e. }(\underline{\Phi}^i_\nu)_{\deleted{\unsure{\bar{x}}}\bar{x}'}=0\quad\text{pour tte gén. }x'\ne x\text{ dans }X_0\ \deleted{\unsure{en}}.
\]
Condition satisfaite en particulier si $x$ est un point maximal de l'ensemble $Z$ des pts de $X_0$ en lesquels $f$ n'est pas lisse, \deleted{\unsure{et lorsque plus il en est|et lorsque plus il a} (4.4.3) \unsure{une comp.}\ \unsure{conn.}\ point de $X$ \unsure{dominant et contenant} $x$}.

L'inclusion
\[
V_{\bar{\unsure{\eta|y}}}\longleftarrow\bar V
\]
définit donc un isomorphisme
\[
\tag*{(6.1.2)}
\mathbb{R}\unsure{\Gamma|\mathscr{T}}(\bar V,\Lambda_\nu)\xrightarrow{\sim}\mathbb{R}\Gamma(V_{\bar\eta},\Lambda_\nu)
\]
i.e. des isomorphismes
\[
\tag*{(6.1.3)}H^i(\bar V,\Lambda_\nu)\simeq H^i(V_{\bar\eta},\Lambda_\nu),
\]
On peut \unsure{en sus} \unsure{compter|considérer} à opérateurs $\unsure{\pi|\Pi}$, (6.1.2) \deleted{\unsure{donne}} s'écrit
\[
\tag*{(6.1.4)}\mathbb{R}(\Gamma,\pi)(\bar V,\Lambda_\nu)\xrightarrow{\sim}\mathbb{R}(\Gamma,\pi)(V_{\bar\eta},\Lambda_\nu\unsure{,|)}.
\]
Donc, au point de vue de leur structure \unsure{chomologique|cohomologique}, la situation des 5 schémas (4.2) se réduit aux quatre schémas $V_s\simeq V_{\bar s},V,V_\eta,V_{\bar\eta}$.

(6.2) L'inclusion $V_{\bar s}\longrightarrow\bar V$ donne donc un homomorphisme \inserted{${\unsure{K|\kappa}}^{\unsure{\bullet|\prime}}$}
\inserted{($(L^\bullet)^{\mathrm{triv}}$, où $L^\bullet=\mathbb{R}\Gamma(V_s,\Lambda_\nu)$)}
\[
\tag*{(6.2.1)}
\boxed{u=(\kappa_\pi^{\unsure{*|\bullet}})^{\unsure{*|\bullet}}:\mathbb{R}(\Gamma,\pi)(V_{\bar\eta},\Lambda_\nu)\longrightarrow\mathbb{R}(\Gamma,\pi)(V_{\bar s},\Lambda_\nu),}
\]
\margin{compte tenu que $\pi$ opère trivialement sur $V_{\bar s}=V_s$}
qui \deleted{par le} s'insère dans le triangle distingué
\diagram{\[
\tag*{(6.2.2)}
\begin{tikzcd}[ampersand replacement=\&,column sep=small]
\& {\mathbb{R}(\Gamma,\pi)(V_{\bar s},\Lambda_\nu)\simeq\mathbb{R}\Gamma(V_s,\Lambda_\nu)^{\mathrm{triv}}=L^{\bullet\,\mathrm{triv}}}\arrow[dl] \& \\
{R^\bullet=\mathbb{R}(\Gamma,\pi)(\bar V,(\Lambda_\nu)_{V_{\bar\eta},\bar V})}\arrow[rr] \& \& {\mathbb{R}(\Gamma,\pi)(V_{\bar\eta},\Lambda_\nu)=K^\bullet}\arrow[ul]
\end{tikzcd}
\]}
